% C7_brief.tex -- ONE source for the English and Romanian case brief of Case 7 (Lecture 7).
\input{numbers_\lang}
\renewcommand{\emqacase}{\T{Research case 7 (Lecture 7)}{Studiul de caz 7 (cursul 7)}}
\begin{document}
\emqatitle{\T{Case 7: An explainable early-warning system for price spikes}%
                {Studiul de caz 7: Un sistem explicabil de avertizare timpurie pentru vîrfurile de preț}}%
          {\T{Based on}{Studiu bazat pe lucrarea} Andrei, Pele, Cramer, Amza \& Diaconu (2026), \emph{Explainable AI-based early warning system for electricity price spikes: evidence from the Romanian market},
           PICBE 20(1), 4820--4836, \href{https://doi.org/10.2478/picbe-2026-0352}{doi:10.2478/picbe-2026-0352}}

\begin{scenario}
\T{A large industrial consumer in Romania (a steel or cement plant) can shift part of its load out of an hour if it knows a day in advance; a supplier can hedge or call demand response; a battery operator can keep energy for the evening.
Each of them asks the same question every morning, before the day-ahead auction closes at noon: \textbf{will tomorrow have a price spike, in which hours, and why?}
The head of energy procurement must decide \textbf{when to act on an alert} --- acting costs money (lost production, a hedge premium), not acting in a spike costs more.
And the board wants to know \textbf{why} the model raises an alarm before trusting it.}%
{Un mare consumator industrial din România (o oțelărie sau o fabrică de ciment) își poate muta o parte din consum dintr-o anumită oră în alta dacă este avertizat cu o zi înainte asupra unui vîrf de preț; un furnizor își poate acoperi riscul sau poate apela la răspunsul cererii; un operator de baterii poate păstra energia stocată pentru seară.
Toți își pun în fiecare dimineață aceeași întrebare, înainte de închiderea, la prînz, a licitației pentru ziua următoare: \textbf{va avea ziua de mîine un vîrf de preț, în ce ore și de ce?}
Directorul de achiziții de energie trebuie să decidă \textbf{cînd să reacționeze la o alertă} --- reacția are un cost (producție pierdută, prima unei acoperiri), dar lipsa de reacție într-un vîrf de preț costă mai mult.
Iar consiliul de administrație dorește să știe \textbf{de ce} modelul emite alerta, înainte de a-i acorda încredere.}
\end{scenario}

\section{\T{What the paper did}{Lucrarea de referință}}
\textbf{\T{Question.}{Întrebarea de cercetare.}} \T{Can explainable artificial intelligence (XAI) identify the drivers of price spikes in the Romanian day-ahead market; is the residual load forecast a primary predictor; and can an operational early-warning system be built from interpretable balance indicators (p.~4821)?}%
{Poate inteligența artificială explicabilă (XAI, \emph{explainable artificial intelligence}) să identifice factorii vîrfurilor de preț de pe piața pentru ziua următoare din România? Este prognoza consumului rezidual una dintre principalele variabile explicative? Se poate construi un sistem operațional de avertizare timpurie pe baza unor indicatori de echilibru ușor de interpretat (p.~4821)?}

\textbf{\T{Data.}{Datele.}} \T{5,000 consecutive hours of the OPCOM day-ahead market, March--September 2024; ENTSO-E generation, load and day-ahead forecasts, oil, gas and carbon prices, temperature in 10 cities; 70 features; chronological split 80/10/10 (p.~4824--4826).}%
{5.000 de ore consecutive de pe piața pentru ziua următoare administrată de OPCOM, martie--septembrie 2024; producția, consumul și prognozele ENTSO-E pentru ziua următoare, prețurile petrolului, gazelor naturale și certificatelor de emisii de CO$_2$, temperaturile din 10 orașe; 70 de variabile; eșantion împărțit cronologic în proporția 80/10/10 (p.~4824--4826).}

\textbf{\T{Method.}{Metoda.}} \T{Random Forest and XGBoost (eXtreme Gradient Boosting) forecasts explained with SHAP (SHapley Additive exPlanations) and LIME (Local Interpretable Model-agnostic Explanations);
a warning score that combines nine indicator thresholds, weighted by their validation F1-score, with an alert threshold $W^*=0.477$ chosen on the validation set; a high-price event is a price above the 90th percentile, 149.99 EUR/MWh (p.~4827, 4832).}%
{Prognoze realizate cu Random Forest (păduri aleatoare) și XGBoost (\emph{eXtreme Gradient Boosting}, arbori de decizie construiți secvențial), explicate cu metodele SHAP (\emph{SHapley Additive exPlanations}) și LIME (\emph{Local Interpretable Model-agnostic Explanations});
un scor de avertizare care combină praguri stabilite pentru nouă indicatori, ponderați cu scorul F1 obținut pe eșantionul de validare, cu un prag de alertă $W^*=0{,}477$ ales tot pe eșantionul de validare; un eveniment de preț ridicat este definit ca un preț peste percentila 90, adică peste 149,99 EUR/MWh (p.~4827, 4832).}

\textbf{\T{Main findings of the paper.}{Principalele rezultate ale lucrării.}}
\begin{itemize}
\item \T{XGBoost explains 67.6\,\% of price variance on the test set ($R^2$ = 0.676, MAE 24.58 EUR/MWh) (p.~4828, Table~3); five-fold time-series cross-validation adds regime-specific evidence: mean $R^2$ 0.542, with the June--July heatwave fold at 0.357 (Table~4).}%
{Pe eșantionul de test, XGBoost explică 67,6\,\% din varianța prețului ($R^2 = 0{,}676$, MAE 24,58 EUR/MWh) (p.~4828, Tabelul~3); validarea încrucișată pentru serii de timp, cu cinci subeșantioane, aduce informații despre regimuri: $R^2$ mediu este 0,542, iar pentru subeșantionul care cuprinde valul de căldură din iunie--iulie valoarea este 0,357 (Tabelul~4).}
\item \T{The most influential features are the previous hour's price (mean $|$SHAP$|$ = 14.30) and the residual load forecast (10.13) (Table~5, p.~4830).}%
{Variabilele cu cea mai mare influență sînt prețul din ora anterioară (media $|$SHAP$|$ = 14,30) și prognoza consumului rezidual (10,13) (Tabelul~5, p.~4830).}
\item \T{The warning system reaches precision 0.470, recall 0.722 and F1 0.569 on the test set (Table~7, p.~4832); the residual-load indicator alone has a validation F1 of 0.792 (Table~8).
The paper also reports a validation F1 of 0.800 and documents the spike frequencies of the two periods (24.0\,\% in validation, 10.8\,\% in test) (p.~4833).}%
{Pe eșantionul de test, sistemul de avertizare obține o precizie de 0,470, o sensibilitate de 0,722 și un scor F1 de 0,569 (Tabelul~7, p.~4832); indicatorul consumului rezidual, luat singur, are un scor F1 de 0,792 pe eșantionul de validare (Tabelul~8).
Lucrarea raportează și un scor F1 de 0,800 pe eșantionul de validare și arată frecvența vîrfurilor în cele două perioade (24,0\,\% în validare, 10,8\,\% în test) (p.~4833).}
\item \T{Missing a spike is ``far more costly'' than a false alarm, so the system favours recall (p.~4832); ``interpretable balance signals retain their directional value'', complementing the price forecasts (p.~4834).}%
{Ratarea unui vîrf costă mult mai mult decît o alarmă falsă, de aceea sistemul acordă prioritate sensibilității (p.~4832); semnalele de echilibru ușor de interpretat își păstrează valoarea ca indicatori de direcție și completează prognozele de preț (p.~4834).}
\end{itemize}
{\footnotesize\color{FooterGray}\T{Related work of our group for Lecture 7: probabilistic price forecasts with distributional random forests (Andrei, Cepoi, Dumitrescu, Lessmann, Chiriac (Luchian), Pele \& Petukhina, 2026, SSRN 6872381, \href{https://doi.org/10.2139/ssrn.6872381}{doi:10.2139/ssrn.6872381}; Case~6); pretrained foundation models with conformal calibration for tail-risk forecasts in CEE markets (Jheng, Tak, Găman, Mazurencu-Marinescu-Pele \& Pele, 2026, ECECSR 60(2), 267--290, \href{https://doi.org/10.24818/18423264/60.2.26.14}{doi:10.24818/18423264/60.2.26.14}).}%
{Alte lucrări ale grupului nostru relevante pentru cursul 7: prognoze probabilistice ale prețului cu păduri aleatoare distribuționale (Andrei, Cepoi, Dumitrescu, Lessmann, Chiriac (Luchian), Pele și Petukhina, 2026, SSRN 6872381, \href{https://doi.org/10.2139/ssrn.6872381}{doi:10.2139/ssrn.6872381}; studiul de caz~6); modele fundamentale preantrenate, cu calibrare conformală, pentru prognoza riscului extrem pe piețele din Europa Centrală și de Est (Jheng, Tak, Găman, Mazurencu-Marinescu-Pele și Pele, 2026, ECECSR 60(2), 267--290, \href{https://doi.org/10.24818/18423264/60.2.26.14}{doi:10.24818/18423264/60.2.26.14}).}}

\section{\T{The method, step by step}{Metoda, pas cu pas}}
\begin{intuition}
\T{Spikes happen when the system is tight: demand is high, the sun has set, the wind is weak, and yesterday was already expensive.
A machine-learning model learns these combinations from history; SHAP then splits each alarm into the contribution of every input, so a manager can see \emph{why}.
The indicator score is a simpler, fully transparent version: a checklist of warning lights, each with a threshold.}%
{Vîrfurile de preț apar atunci cînd sistemul este sub presiune: cererea este mare, soarele a apus, vîntul este slab, iar prețul din ziua precedentă a fost deja ridicat.
Un model de învățare automată deduce aceste combinații din datele istorice; metoda SHAP descompune apoi fiecare alertă în contribuțiile variabilelor de intrare, astfel încît managerul înțelege \emph{de ce} a fost emisă.
Scorul bazat pe indicatori este o variantă mai simplă și complet transparentă: o listă de semnale de avertizare, fiecare cu propriul prag.}
\end{intuition}
\T{A spike in hour $t$ is $S_t=\mathbb 1\{p_t>\theta\}$, with $\theta$ the 90th percentile of training prices. The residual load is $RL_t=L_t-\mathit{Solar}_t-\mathit{Wind}_t$ (paper eq.~1, with day-ahead forecasts).
The classifier estimates $\hat\pi_t=\Pr(S_t=1\mid\mathbf x_t)$ with gradient-boosted trees ($\hat f(\mathbf x)=\sum_{k=1}^{K}f_k(\mathbf x)$, each tree $f_k$ correcting the previous ones).
SHAP writes each prediction as a sum of contributions, $\hat f(\mathbf x_t)=\phi_0+\sum_j\phi_{j,t}$, where $\phi_{j,t}$ is the average change in the prediction when feature $j$ is added, over all orders of adding the features (the Shapley value from game theory).
The indicator score is $W_t=\sum_m w_m\,\mathbb 1\{x_{m,t}\ge\theta_m\}$, with thresholds $\theta_m$ and weights $w_m\propto F1_m$; an alert is raised when $W_t\ge W^*$.
Quality: precision $=TP/(TP+FP)$ (share of alerts that were spikes), recall $=TP/(TP+FN)$ (share of spikes that were caught), $F1=2\cdot$precision$\cdot$recall$/($precision$+$recall$)$, with $TP$, $FP$, $FN$ the true positives, false positives and false negatives.
\textbf{Decision rule:} if acting costs $C$ per MWh and an unhedged spike costs $L$ per MWh, act whenever $\hat\pi_t>C/L$ --- for an operating decision, the company's own cost ratio sets the threshold.}%
{Un vîrf în ora $t$ este definit prin $S_t=\mathbb 1\{p_t>\theta\}$, unde $\theta$ este percentila 90 a prețurilor din eșantionul de antrenare. Consumul rezidual este $RL_t=L_t-\mathit{Solar}_t-\mathit{Eolian}_t$ (ec.~1 din lucrare, cu prognoze pentru ziua următoare).
Clasificatorul estimează $\hat\pi_t=\Pr(S_t=1\mid\mathbf x_t)$ cu arbori construiți secvențial ($\hat f(\mathbf x)=\sum_{k=1}^{K}f_k(\mathbf x)$, fiecare arbore $f_k$ corectînd erorile celor anteriori).
Metoda SHAP descompune fiecare prognoză într-o sumă de contribuții, $\hat f(\mathbf x_t)=\phi_0+\sum_j\phi_{j,t}$, unde $\phi_{j,t}$ este variația medie a prognozei la adăugarea variabilei $j$, calculată pe toate ordinile posibile de adăugare a variabilelor (valoarea Shapley din teoria jocurilor).
Scorul bazat pe indicatori este $W_t=\sum_m w_m\,\mathbb 1\{x_{m,t}\ge\theta_m\}$, cu pragurile $\theta_m$ și ponderile $w_m\propto F1_m$; alerta se emite atunci cînd $W_t\ge W^*$.
Măsurile de calitate: precizia $=TP/(TP+FP)$ (proporția alertelor care au fost urmate de vîrfuri), sensibilitatea $=TP/(TP+FN)$ (proporția vîrfurilor semnalate), $F1=2\cdot$precizie$\cdot$sensibilitate$/($precizie$+$sensibilitate$)$, unde $TP$, $FP$ și $FN$ reprezintă, respectiv, numărul rezultatelor pozitive adevărate, al celor pozitive false și al celor negative false.
\textbf{Regula de decizie:} dacă reacția costă $C$ pe MWh, iar un vîrf neacoperit costă $L$ pe MWh, reacționați de fiecare dată cînd $\hat\pi_t>C/L$ --- pentru o decizie operațională, pragul este stabilit de raportul dintre costurile proprii.}
\begin{decision}
\T{Use the model to rank hours by spike probability and to explain them; set the action threshold from your own costs ($C/L$); recalibrate thresholds when the market regime changes.}%
{Folosiți modelul pentru a ordona orele după probabilitatea unui vîrf și pentru a explica aceste probabilități; stabiliți pragul de reacție pe baza propriilor costuri ($C/L$); recalibrați pragurile atunci cînd se schimbă regimul pieței.}
\end{decision}

\section{\T{Replicate it yourself}{Replicare pe datele cursului}}
\T{Notebook \texttt{C7\_notebook\_EN.ipynb}: hourly prices, load, solar, wind, hydro, gas and net imports (energy-charts.info). Training 2023--2024, validation January--June 2025, test July 2025 -- October 2026 (\NnTest\ hours).
Spike threshold: \Nthr\ EUR/MWh (P90 of training; paper: 149.99). Spike share: \NspikeTrain\,\% in training, \NspikeVal\,\% in validation, \NspikeTest\,\% in the test period.
Classroom implementation: our data contain no day-ahead forecasts, so the main models use only the information available at gate closure (12:00 CET on $d-1$): prices up to $d-1$ and the realised load, solar, wind, hydro, gas and imports of $d-2$ (day $d-1$ is not yet metered); a second XGBoost classifier with the actual fundamentals of the delivery hour (perfect foresight) is shown for comparison only. No weather or fuels; a longer sample with a different regime.}%
{Notebook-ul \texttt{C7\_notebook\_RO.ipynb} folosește prețurile orare, consumul, producția solară, eoliană, hidro și pe gaze, precum și importurile nete (energy-charts.info). Eșantionul de antrenare acoperă 2023--2024, cel de validare ianuarie--iunie 2025, iar cel de test iulie 2025 -- octombrie 2026 (\NnTest\ ore).
Pragul de vîrf este de \Nthr\ EUR/MWh (percentila 90 din eșantionul de antrenare; în lucrare: 149,99). Ponderea orelor cu vîrfuri: \NspikeTrain\,\% în antrenare, \NspikeVal\,\% în validare, \NspikeTest\,\% în perioada de test.
Față de lucrare, implementarea de la curs ține seama de faptul că datele noastre nu conțin prognoze pentru ziua următoare: modelele principale folosesc numai informația disponibilă la închiderea porții (ora 12:00 CET în ziua $d-1$), adică prețurile pînă în ziua $d-1$, precum și consumul, producția solară, eoliană, hidro, pe gaze și importurile realizate în ziua $d-2$ (ziua $d-1$ nu este încă măsurată); un al doilea clasificator XGBoost, cu valorile efective ale variabilelor fundamentale din ora de livrare (previziune perfectă), este prezentat numai pentru comparație. Implementarea nu include variabile meteorologice sau prețurile combustibililor și acoperă un eșantion mai lung, cu un alt regim de piață.}

\begin{figure}[H]
\begin{minipage}[t]{0.52\linewidth}\centering
\includegraphics[width=\linewidth]{f1_spikes}
\caption{\T{Daily maximum price and the spike threshold.}{Prețul orar maxim al fiecărei zile și pragul de vîrf.}}
\src{energy-charts.info; 12.2022--10.2026.}
\end{minipage}\hfill
\begin{minipage}[t]{0.46\linewidth}\centering
\includegraphics[width=\linewidth]{f2_shap_importance}
\caption{\T{Top-10 drivers of spike alerts (SHAP).}{Primii 10 factori ai alertelor de vîrf de preț (SHAP).}}
\src{\T{our XGBoost classifier; test period.}{clasificatorul nostru XGBoost; perioada de test.}}
\end{minipage}
\end{figure}

\begin{table}[H]\centering\small
\caption{\T{Early-warning performance on the test period (July 2025 -- October 2026). Best value in bold.}{Performanța avertizării timpurii în perioada de test (iulie 2025 -- octombrie 2026). Cea mai bună valoare este marcată cu aldine.}}
\tabin{warning}
\src{\T{our computation. Paper (test set, 2024): warning system precision 0.470, recall 0.722, F1 0.569 (Table~7).}{calcule proprii. Lucrarea (eșantionul de test, 2024): sistemul de avertizare are precizia 0,470, sensibilitatea 0,722 și F1 0,569 (Tabelul~7).}}
\end{table}

\T{\textbf{What our data show.} With gate-closure information only, the top three SHAP features are \Ntop, \Nsecond\ and \Nthird: past prices dominate. Our price lag is the same hour of day $d-1$, the latest price known when bidding for day $d$; the paper's other main driver, the residual load \emph{forecast} (Table~5), is not in our data and enters our model only through its $d-2$ value.
The ML classifier separates spike and non-spike hours well (AUC \NxAUC); the indicator score is weaker (AUC \NwAUC, F1 \NwFone). The actual fundamentals of the hour (perfect foresight) do not improve the test results (AUC \NoAUC, F1 \NoFone), a sign that the link between fundamentals and spikes learnt in 2023--2024 changed in the test period.
\textbf{But a regime shift dominates the test period:} spikes became daily evening events (\NspikeTest\,\% of hours vs \NspikeTrain\,\% in training), so the thresholds tuned on 2025 validation miss many of them (classifier recall \NxRec), and the naive rule ``a spike yesterday at the same hour'' has the highest recall and F1 (\NnFone).
The managerial lesson: when the market regime changes, recalibrate thresholds and keep a persistence benchmark in the monitoring.}%
{\textbf{Ce arată datele.} Numai cu informația de la închiderea porții, primele trei variabile după importanța SHAP sînt \Ntop, \Nsecond\ și \Nthird: domină prețurile anterioare. Decalajul nostru de preț este aceeași oră din ziua $d-1$, cel mai recent preț cunoscut la licitația pentru ziua $d$; celălalt factor principal din lucrare, \emph{prognoza} consumului rezidual (Tabelul~5), nu se află în datele noastre și intră în model numai prin valoarea sa din ziua $d-2$.
Clasificatorul separă bine orele cu vîrfuri de celelalte (AUC \NxAUC); scorul bazat pe indicatori este mai slab (AUC \NwAUC, F1 \NwFone). Valorile efective ale variabilelor fundamentale din ora respectivă (previziune perfectă) nu îmbunătățesc rezultatele din perioada de test (AUC \NoAUC, F1 \NoFone), semn că legătura dintre variabilele fundamentale și vîrfuri, învățată pe 2023--2024, s-a modificat în perioada de test.
\textbf{Perioada de test este însă dominată de o schimbare de regim:} vîrfurile au devenit evenimente zilnice, de seară (\NspikeTest\,\% din ore, față de \NspikeTrain\,\% în antrenare), astfel că pragurile calibrate pe perioada de validare din 2025 ratează multe dintre ele (sensibilitatea clasificatorului este \NxRec), iar regula naivă „vîrf ieri, la aceeași oră” obține cele mai bune valori pentru sensibilitate și F1 (\NnFone).
Concluzia pentru management: atunci cînd se schimbă regimul pieței, recalibrați pragurile și păstrați în monitorizare un reper de persistență.}

\begin{figure}[H]
\begin{minipage}[t]{0.42\linewidth}\centering
\includegraphics[width=\linewidth]{f3_local_explanation}
\caption{\T{Why the alarm on the highest test price (\NpeakTime, \NpeakPrice\ EUR/MWh)?}{De ce a fost emisă alerta pentru cel mai mare preț din perioada de test (\NpeakTime, \NpeakPrice\ EUR/MWh)?}}
\src{\T{our SHAP values.}{valorile noastre SHAP.}}
\end{minipage}\hfill
\begin{minipage}[t]{0.56\linewidth}\centering
\includegraphics[width=\linewidth]{f4_decision}
\caption{\T{Reliability of the spike probability (left) and value of acting by threshold (right).}{Fiabilitatea probabilității de vîrf (stînga) și valoarea reacției în funcție de prag (dreapta).}}
\src{\T{our computation; test period; 1 MW flexible load, $C$ = 30 EUR/MWh.}{calcule proprii; perioada de test; 1 MW de consum flexibil, $C$ = 30 EUR/MWh.}}
\end{minipage}
\end{figure}

\T{\textbf{Decision number.} At the highest test price the model gave a spike probability of \NpeakProb. With an action cost $C$ = 30 EUR/MWh and a loss $L$ = \Nloss\ EUR/MWh (mean spike minus mean non-spike price in training), the cost--loss threshold is $C/L$ = \Npc.
Acting at that threshold with 1 MW of flexible load would have saved about \NvalPc\ thousand EUR over the test period, versus \NvalPstar\ thousand at the F1-optimal threshold --- for an operating decision, the threshold should reflect the company's own costs.}%
{\textbf{Cifra relevantă pentru decizie.} Pentru cel mai mare preț din perioada de test, modelul a estimat o probabilitate de vîrf de \NpeakProb. La un cost al reacției $C$ = 30 EUR/MWh și o pierdere $L$ = \Nloss\ EUR/MWh (diferența dintre prețul mediu în orele cu vîrfuri și cel din celelalte ore, în eșantionul de antrenare), pragul cost--pierdere este $C/L = \Npc$.
Reacționînd la acest prag cu 1 MW de consum flexibil, s-ar fi economisit aproximativ \NvalPc\ mii EUR în perioada de test, față de \NvalPstar\ mii EUR la pragul optim după F1 --- pentru o decizie operațională, pragul trebuie să reflecte costurile proprii ale companiei.}

\section{\T{Discussion questions for managers}{Întrebări pentru discuție}}
\begin{enumerate}
\item \T{In our 2025--2026 test, the persistence rule scores highest. What would you use each tool for, and how would you monitor and recalibrate them?}%
{În perioada noastră de test, 2025--2026, regula de persistență obține cele mai bune scoruri. În ce scop ați folosi fiecare instrument și cum le-ați monitoriza și recalibra?}
\item \T{Your board must approve an AI tool. Is a SHAP explanation enough for them to trust an alarm? What else would you show?}%
{Consiliul de administrație trebuie să aprobe un instrument bazat pe inteligență artificială. Este suficientă o explicație SHAP pentru a avea încredere într-o alertă? Ce alte informații ați prezenta?}
\item \T{What are $C$ and $L$ for your company? How would the alert threshold change for a cement plant and for a battery operator?}%
{Care sînt valorile $C$ și $L$ pentru compania dumneavoastră? Cum s-ar modifica pragul de alertă pentru o fabrică de ciment și pentru un operator de baterii?}
\item \T{How often should thresholds be recalibrated, and who in the company should own that decision?}%
{Cît de des ar trebui recalibrate pragurile și cine din companie ar trebui să răspundă de această decizie?}
\item \T{The spike threshold is a percentile of the past. When spikes become daily, is ``spike'' still the right concept, or should you forecast the whole price distribution (Case~6)?}%
{Pragul de vîrf este o percentilă calculată pe date trecute. Cînd vîrfurile devin zilnice, mai este „vîrful” conceptul potrivit sau ar trebui prognozată întreaga distribuție a prețului (studiul de caz~6)?}
\end{enumerate}

\section{\T{Limitations of the classroom replication}{Limitele replicării}}
\T{(i) Our models and thresholds are fitted once (training 2023--2024, validation January--June 2025) and not re-estimated during the test. (ii) Our main models replace the day-ahead forecasts of load, solar and wind by their values on $d-2$, the last metered day, and use no weather or fuel prices (both pessimistic).
(iii) The spike definition depends on the training percentile, which drifts with the price level. (iv) The cost--loss value assumes the flexible load can always be shifted and ignores the price impact of doing so.}%
{(i) Modelele și pragurile noastre sînt estimate o singură dată (antrenare 2023--2024, validare ianuarie--iunie 2025) și nu sînt reestimate în perioada de test. (ii) Modelele principale înlocuiesc prognozele pentru ziua următoare ale consumului și ale producției solare și eoliene cu valorile din ziua $d-2$, ultima zi măsurată, și nu includ variabile meteorologice sau prețurile combustibililor (ambele alegeri le dezavantajează).
(iii) Definiția vîrfului depinde de percentila calculată pe eșantionul de antrenare, care se modifică odată cu nivelul prețurilor. (iv) Calculul cost--pierdere presupune că, de fiecare dată, consumul flexibil poate fi mutat și ignoră efectul acestei mutări asupra prețului.}

\newpage
\section*{\T{Case 7 at a glance}{Studiul de caz 7 pe scurt}}
\glance{\T{Will tomorrow have a price spike, why, and when should we act?}{Va avea ziua de mîine un vîrf de preț, de ce și cînd trebuie să reacționăm?}}%
{\T{Hourly RO prices and fundamentals 2023--Oct 2026; spike = price above P90 of training.}{Prețuri orare RO și variabile fundamentale, 2023--oct.\ 2026; vîrf = preț peste percentila 90 din antrenare.}}%
{\T{XGBoost + SHAP; F1-weighted indicator score; cost--loss decision rule.}{XGBoost + SHAP; scor bazat pe indicatori ponderați cu F1; regulă de decizie cost--pierdere.}}%
{\T{Gate-closure information: top drivers \Ntop, \Nsecond; AUC \NxAUC; regime shift: naive F1 \NnFone.}{Informația de la închiderea porții: factorii principali sînt \Ntop\ și \Nsecond; AUC \NxAUC; schimbare de regim: F1 naiv \NnFone.}}%
{\T{Act when spike probability $> C/L$ (\Npc); recalibrate thresholds as the market changes.}{Reacționați cînd probabilitatea unui vîrf $> C/L$ (\Npc); recalibrați pragurile cînd se schimbă piața.}}
\begin{center}
\includegraphics[width=0.62\linewidth]{fig_headline_\lang}
\src{\T{our XGBoost classifier on energy-charts.info data; SHAP on the test period, July 2025 -- October 2026.}{clasificatorul nostru XGBoost pe date energy-charts.info; SHAP pe perioada de test, iulie 2025 -- octombrie 2026.}}
\end{center}
\begin{minipage}[t]{0.32\linewidth}\centering{\color{MainBlue}\Huge\bfseries \NxAUC}\\ \small \T{AUC of the spike classifier}{AUC al clasificatorului de vîrfuri}\end{minipage}\hfill
\begin{minipage}[t]{0.32\linewidth}\centering{\color{MainBlue}\Huge\bfseries \NspikeTest\,\%}\\ \small \T{of test hours were spikes (training: \NspikeTrain\,\%)}{din orele de test au fost ore cu vîrfuri (antrenare: \NspikeTrain\,\%)}\end{minipage}\hfill
\begin{minipage}[t]{0.32\linewidth}\centering{\color{IDAred}\Huge\bfseries \NvalPc}\\ \small \T{thousand EUR saved per MW at $C/L$}{mii EUR economisiți pe MW la pragul $C/L$}\end{minipage}
\end{document}
